\documentclass[12pt,a4paper,oneside]{book}
\usepackage[utf8]{vietnam}
\usepackage[utf8]{inputenc}
\usepackage{amssymb}
\usepackage{graphicx}
\usepackage{subcaption}
\usepackage{booktabs}
\usepackage{mathptmx}
\usepackage{amsmath}
\usepackage{amssymb}
\usepackage{amsfonts}
\usepackage{array}
\newcolumntype{P}[1]{>{\centering\arraybackslash}p{#1}}
\newcolumntype{M}[1]{>{\centering\arraybackslash}m{#1}}
\usepackage{fancyhdr}
\usepackage{multirow}
\usepackage{algorithm}
\usepackage{algpseudocode}
\makeatletter
\renewcommand{\ALG@name}{Giải thuật}
\renewcommand{\listalgorithmname}{Danh sách giải thuật}
\makeatother
\def\algbackskip{\hskip-\ALG@thistlm}
\makeatother
\usepackage{caption}
\captionsetup[algorithm]{labelsep=colon}

% Set line spacing
\linespread{1.25} 
% Set paragraph spacing
\setlength{\parskip}{\baselineskip}%
\usepackage[skip=10pt plus1pt, indent=15pt]{parskip}

% Colors and hyperlinks
\usepackage{float}
\usepackage{colortbl}
\usepackage[usenames,dvipsnames]{xcolor}
\definecolor{cobalt}{rgb}{0.0, 0.28, 0.67}
\definecolor{carmine}{rgb}{0.59, 0.0, 0.09}
\usepackage[unicode, colorlinks=true, urlcolor=carmine, citecolor=cobalt, linkcolor=carmine, pdfborder={0 0 0}]{hyperref}

% Thesis style
\usepackage{styles/thesis}
\graphicspath{{figs/}}
\setcounter{secnumdepth}{3}

% Skip blank page at beginning
\usepackage{atbegshi}
\AtBeginDocument{\AtBeginShipoutNext{\AtBeginShipoutDiscard}}

% For code highlighting
\usepackage{listings}
\lstset{
    basicstyle=\ttfamily\small,
    breaklines=true,
    frame=single,
    numbers=left,
    numberstyle=\tiny,
    tabsize=2,
    captionpos=b
}

% For long tables
\usepackage{longtable}
\usepackage{enumitem}


% === THÔNG TIN ĐỒ ÁN ===
% TODO: Cập nhật thông tin bên dưới cho đúng
\ctname{THIẾT KẾ VÀ TRIỂN KHAI HỆ THỐNG ZERO TRUST SOFTWARE SUPPLY CHAIN TRÊN KUBERNETES THEO CHUẨN SLSA LEVEL 3}
\cstuname{Họ và Tên Sinh Viên}
\csCouncil{Công nghệ thông tin}
\csSupervise{TS. Tên Giảng Viên}
\cttime{2026}
% \csdeptname{KHOA CÔNG NGHỆ THÔNG TIN}

\thesislayout

\begin{document}
% Trang bìa
\coverpage

\frontmatter

% Lời cam đoan
\begin{declaration}
Tôi cam đoan đây là công trình do tôi tự thực hiện dưới sự hướng dẫn của TS. Tên Giảng Viên. Các nội dung nghiên cứu, số liệu và kết quả là trung thực. Các số liệu, công trình sử dụng của tác giả khác đều được trích dẫn nguồn gốc rõ ràng. Tất cả phần mềm, mã nguồn được sử dụng trong đồ án này đều có giấy phép mã nguồn mở.

Nếu phát hiện có bất kỳ sự gian lận nào, tôi xin chịu hoàn toàn trách nhiệm.

\begin{flushright}
{Họ và Tên Sinh Viên}
\end{flushright}
\end{declaration}

\newpage

% Lời cảm ơn
\begin{acknowledgments}
% TODO: Viết lời cảm ơn
Lời cảm ơn sẽ được bổ sung sau.
\end{acknowledgments}

% Tóm tắt
\begin{abstract}
% TODO: Viết tóm tắt sau khi hoàn thành đồ án
Đồ án nghiên cứu và triển khai hệ thống Zero Trust Software Supply Chain cho Kubernetes theo chuẩn SLSA Level 3, nhằm đảm bảo an toàn cho quy trình build và triển khai container image. Hệ thống bao gồm pipeline CI/CD tự động tạo provenance và ký số container image, quản lý SBOM, và Kubernetes Admission Controller kiểm soát việc triển khai container theo mô hình Zero Trust. Kết quả thực nghiệm cho thấy hệ thống có khả năng ngăn chặn hiệu quả các kịch bản tấn công giả mạo và chèn mã độc trong chuỗi cung ứng phần mềm.

\textbf{Từ khóa:} Zero Trust, Software Supply Chain, SLSA, Kubernetes, Container Security, Sigstore, Provenance, SBOM.
\end{abstract}

% Mục lục
\tableofcontents
% Danh mục bảng
\listoftables
% Danh mục hình
\listoffigures

\mainmatter
\fancyhead{}
\renewcommand{\footrulewidth}{0.4pt}
\pagestyle{fancy} 
\renewcommand{\chaptermark}[1]{\markboth{#1}{#1}}
\fancyhead[R]{\chaptername\ \thechapter\ --\ \leftmark}

% === CÁC CHƯƠNG ===
\chapter{Tổng quan đề tài}\label{ch:introduction}

\section{Bối cảnh và lý do chọn đề tài}

Trong bối cảnh công nghệ thông tin hiện đại, an ninh chuỗi cung ứng phần mềm (Software Supply Chain Security) đã trở thành một trong những ưu tiên chiến lược cấp bách nhất. Các cuộc tấn công chấn động như SolarWinds (2020), Codecov (2021), Log4Shell (2021) và xz-utils (2024) đã phơi bày một thực tế: tin tặc không chỉ nhắm vào ứng dụng trực tiếp mà đang khai thác các điểm yếu trong hệ thống hạ tầng xây dựng và các thư viện mã nguồn mở phụ thuộc.

Mô hình bảo mật truyền thống -- chỉ bảo vệ vòng ngoài (perimeter security) -- đã không còn đủ khả năng đối phó với các vector tấn công tinh vi nhắm vào toàn bộ vòng đời phát triển phần mềm. Mô hình Zero Trust với nguyên tắc ``Never trust, always verify'' cung cấp một cách tiếp cận phù hợp hơn, yêu cầu xác minh mọi thành phần tại mọi giai đoạn.

Khung tiêu chuẩn SLSA (Supply-chain Levels for Software Artifacts) -- được phát triển bởi OpenSSF dựa trên các thực tiễn bảo mật của Google -- cung cấp một mô hình trưởng thành với các cấp độ bảo mật tăng dần, giúp tổ chức đánh giá và cải thiện tính toàn vẹn chuỗi cung ứng một cách có hệ thống.

Kubernetes, với vai trò là nền tảng orchestration container phổ biến nhất hiện nay, cung cấp cơ chế Admission Controller cho phép kiểm soát chặt chẽ việc triển khai container -- tạo điều kiện lý tưởng để hiện thực hóa mô hình Zero Trust tại điểm cuối của chuỗi cung ứng.

% TODO: Bổ sung thêm dẫn chứng và trích dẫn

\section{Mục tiêu nghiên cứu}

Các mục tiêu chính của đồ án bao gồm:

\begin{itemize}
    \item Nghiên cứu các công nghệ và tiêu chuẩn liên quan đến DevSecOps, Zero Trust Architecture và Software Supply Chain Security.
    \item Xây dựng pipeline CI/CD theo chuẩn SLSA Level 3, đảm bảo tính toàn vẹn và truy vết nguồn gốc của container image.
    \item Áp dụng cơ chế ký số container image và tạo các bằng chứng xác thực (provenance, attestation).
    \item Tạo và quản lý SBOM (Software Bill of Materials), thực hiện quét và đánh giá lỗ hổng bảo mật.
    \item Triển khai Kubernetes Admission Controller để kiểm soát việc triển khai container theo mô hình Zero Trust.
    \item Thực nghiệm, mô phỏng các kịch bản tấn công và đánh giá hiệu quả của giải pháp.
\end{itemize}

\section{Đối tượng và phạm vi nghiên cứu}

\subsection{Đối tượng nghiên cứu}
\begin{itemize}
    \item Quy trình CI/CD an toàn theo chuẩn SLSA Level 3
    \item Cơ chế ký số và xác thực container image (Sigstore/Cosign)
    \item Kubernetes Admission Controller và Policy Engine
    \item Software Bill of Materials (SBOM) và quét lỗ hổng bảo mật
\end{itemize}

\subsection{Phạm vi nghiên cứu}
\begin{itemize}
    \item \textbf{Nền tảng core:} Kubernetes (triển khai local cluster hoặc cloud)
    \item \textbf{CI/CD:} GitHub Actions hoặc Tekton Pipelines (tùy theo đánh giá khả thi)
    \item \textbf{Signing \& Attestation:} Sigstore ecosystem (Cosign, Fulcio, Rekor)
    \item \textbf{Policy Engine:} Kyverno hoặc tương đương
    \item \textbf{SBOM:} Syft hoặc Trivy (chuẩn SPDX/CycloneDX)
\end{itemize}

\subsection{Giới hạn}
\begin{itemize}
    \item Đồ án tập trung vào khía cạnh bảo mật chuỗi cung ứng, không đi sâu vào phát triển ứng dụng.
    \item Sample application chỉ mang tính minh họa, ưu tiên tính đơn giản.
    \item Không bao gồm: runtime security monitoring, network policies, hermetic builds (SLSA L4).
\end{itemize}

\section{Bố cục đồ án}

Đồ án gồm các phần như sau:
\begin{itemize}
    \item \textbf{Chương \ref{ch:introduction}}: Tổng quan đề tài -- Trình bày bối cảnh, mục tiêu, phạm vi nghiên cứu.
    \item \textbf{Chương \ref{ch:background}}: Cơ sở lý thuyết -- Trình bày nền tảng kiến thức về Supply Chain Security, Zero Trust, SLSA, Kubernetes.
    \item \textbf{Chương \ref{ch:design}}: Phân tích và Thiết kế hệ thống -- Kiến trúc tổng thể và thiết kế chi tiết từng giai đoạn.
    \item \textbf{Chương \ref{ch:implementation}}: Triển khai và Đánh giá -- Cài đặt, cấu hình, kịch bản kiểm thử, đánh giá kết quả.
    \item \textbf{Chương \ref{ch:conclusion}}: Kết luận và Hướng phát triển.
\end{itemize}

\chapter{Cơ sở lý thuyết}\label{ch:background}

\section{Software Supply Chain và các rủi ro bảo mật}

\subsection{Định nghĩa chuỗi cung ứng phần mềm}
% TODO: Định nghĩa chuỗi cung ứng phần mềm
% Source Code → Dependencies → Build Process → Artifacts → Registry → Deployment

\subsection{Các vector tấn công phổ biến}
% TODO: Phân tích các vector tấn công ở từng giai đoạn
% - Source: unauthorized commit, compromised repo
% - Build: compromised build process, cache poisoning
% - Dependency: dependency confusion, typosquatting
% - Distribution: registry poisoning, MITM

\subsection{Các cuộc tấn công tiêu biểu}
% TODO: Case studies
% - SolarWinds (2020): build process compromise
% - Codecov (2021): CI script modification
% - Log4Shell (2021): dependency vulnerability
% - xz-utils (2024): long-term social engineering

\section{Mô hình bảo mật Zero Trust}

\subsection{Khái niệm Zero Trust}
% TODO: Giới thiệu Zero Trust Architecture
% - Nguyên tắc "Never trust, always verify"
% - NIST SP 800-207

\subsection{Áp dụng Zero Trust vào SDLC}
% TODO: Cách áp dụng Zero Trust vào quy trình phát triển phần mềm
% - Không tin tưởng bất kỳ code commit nào chưa verified
% - Không tin tưởng build runner chưa xác thực
% - Không tin tưởng container image chưa có chữ ký hợp lệ
% - Verify at every stage

\section{Khung tiêu chuẩn SLSA}

\subsection{Tổng quan về SLSA}
% TODO: Giới thiệu SLSA
% - Nguồn gốc: Google Binary Authorization for Borg
% - OpenSSF project
% - Maturity model, không phải tool

\subsection{Build Track -- Các cấp độ bảo mật}
% TODO: Bảng phân tích Build L0 → L3
% L0: No guarantees
% L1: Automated build + provenance documentation
% L2: Hosted build platform + signed provenance  
% L3: Isolated/ephemeral build + unforgeable provenance

\subsection{Source Track}
% TODO: Source L1 → L4 (theo draft v1.2)
% Tập trung giải thích L1-L3 (phần liên quan đến đồ án)

\subsection{Threat Model theo SLSA}
% TODO: Phân tích mô hình đe dọa
% - Source threats
% - Build threats  
% - Packaging & Distribution threats
% Ánh xạ với các cấp độ SLSA

\subsection{Mô hình xây dựng và Provenance Schema}
% TODO: 
% - Build Model: Control Plane vs Build Environment
% - In-toto attestation format
% - Provenance schema: buildDefinition, runDetails
% - externalParameters, resolvedDependencies

\section{Kubernetes và Container Security}

\subsection{Kiến trúc Kubernetes}
% TODO: Tổng quan K8s
% - Control Plane: API Server, etcd, Scheduler, Controller Manager
% - Worker Nodes: kubelet, container runtime
% - Pods, Deployments, Services

\subsection{Admission Controllers trong Kubernetes}
% TODO: Giải thích cơ chế Admission Controller
% - Validating Webhooks
% - Mutating Webhooks
% - Luồng xử lý: API Request → Authentication → Authorization → Admission → etcd

\subsection{OCI Image Format và Image Security}
% TODO:
% - OCI image spec
% - Image digest vs tag (tại sao tag không an toàn)
% - Image signing concepts

\section{Hệ sinh thái công cụ liên quan}

\subsection{Sigstore -- Hệ sinh thái ký số}
% TODO: Giới thiệu Sigstore
% - Cosign: ký và xác thực container image
% - Fulcio: certificate authority (keyless signing via OIDC)
% - Rekor: transparency log
% - Keyless signing workflow

\subsection{SBOM -- Software Bill of Materials}
% TODO:
% - Định nghĩa SBOM
% - Các chuẩn: SPDX, CycloneDX
% - Công cụ: Syft, Trivy
% - Vai trò trong supply chain security

\subsection{Policy Engines}
% TODO: So sánh các policy engine
% - Kyverno: YAML-native, dễ sử dụng, tích hợp Cosign
% - OPA Gatekeeper: Rego language, mạnh nhưng phức tạp
% - Sigstore Policy Controller
% Giải thích tại sao chọn phương án linh hoạt

\chapter{Phân tích và Thiết kế hệ thống}\label{ch:design}

\section{Kiến trúc tổng thể}

% TODO: Vẽ sơ đồ kiến trúc tổng thể (tool-agnostic)
% Developer → Source Control → CI/CD Build → Sign & Attest → Registry → Policy Gate → K8s Runtime

Kiến trúc hệ thống được thiết kế theo mô hình phân lớp (layered architecture), trong đó mỗi lớp có thể thay thế công cụ cụ thể mà không ảnh hưởng đến các lớp còn lại. Kubernetes đóng vai trò là nền tảng core xuyên suốt hệ thống.

% TODO: Bảng ánh xạ các lớp kiến trúc với công cụ
% | Layer | Primary Option | Fallback Option |
% | CI/CD | GitHub Actions | Tekton Pipelines |
% | Build | Docker build / Kaniko | Buildah |
% | Sign  | Cosign (keyless) | Cosign (key-based) |
% | Provenance | slsa-github-generator / Tekton Chains | ... |
% | SBOM | Syft | Trivy |
% | Registry | GHCR | Harbor |
% | Policy | Kyverno | Sigstore Policy Controller |

\section{Thiết kế giai đoạn Source -- Quản lý mã nguồn}

% TODO: Thiết kế chi tiết
% - Branch protection rules (chặn force push, chặn xóa nhánh)
% - Commit signing (GPG/SSH)
% - Bắt buộc code review (ít nhất 1 reviewer)
% - Status checks phải pass trước khi merge
% - Secret scanning (TruffleHog / Gitleaks)
% - SAST scanning (Semgrep)

\section{Thiết kế giai đoạn Build -- Đáp ứng SLSA Level 3}

\subsection{Môi trường Build phù du và cô lập}
% TODO: Giải thích yêu cầu ephemeral + isolated
% Phương án A (GitHub Actions):
% - Mỗi workflow run chạy trên runner riêng biệt
% - slsa-github-generator sử dụng reusable workflow để cách ly signing
% Phương án B (Tekton on K8s):
% - Mỗi TaskRun chạy trong Pod riêng, pod bị hủy sau khi hoàn thành
% - Tekton Chains chạy ở Control Plane level

\subsection{Tạo Provenance}
% TODO:
% - In-toto attestation format
% - buildDefinition: buildType, externalParameters (repo URL, commit SHA)
% - runDetails: builder.id, metadata
% - subject: image digest (SHA256)
% Phương án A: slsa-github-generator tự tạo
% Phương án B: Tekton Chains tự observe và tạo

\subsection{Tạo SBOM}
% TODO:
% - Sử dụng Syft hoặc Trivy để quét image
% - Xuất định dạng SPDX hoặc CycloneDX
% - Gắn SBOM vào OCI registry (cosign attach sbom)

\subsection{Ký số -- Digital Signing}
% TODO:
% - Keyless signing workflow qua Sigstore:
%   1. CI system request OIDC token
%   2. Fulcio issue short-lived certificate
%   3. Cosign sign image + provenance + SBOM
%   4. Record to Rekor transparency log
% - Fallback: key-based signing (cosign generate-key-pair)

\section{Thiết kế giai đoạn Deploy -- Zero Trust trên Kubernetes}

\subsection{Kiến trúc Admission Controller}
% TODO: Sơ đồ luồng xử lý
% kubectl apply → API Server → Admission Webhook (Kyverno) → Verify → Allow/Deny

\subsection{Thiết kế Policy Rules}
% TODO: Liệt kê các policy rules
% Policy 1: Chặn image không có chữ ký (unsigned image)
% Policy 2: Chặn image không có SLSA provenance
% Policy 3: Chặn image từ builder không tin cậy (builder.id mismatch)
% Policy 4: Chặn image tham chiếu bằng tag (bắt buộc dùng digest)
% Policy 5: Kiểm tra SBOM đính kèm
% Mô hình: deny-by-default

\section{Thiết kế kịch bản kiểm thử và tấn công mô phỏng}

% TODO: Bảng kịch bản kiểm thử
%
% | # | Kịch bản | Mô tả | Kết quả mong đợi |
% |---|---|---|---|
% | 1 | Happy Path | Image hợp lệ, signed, có provenance | Deploy thành công |
% | 2 | Unsigned Image | Image không có chữ ký | Bị chặn bởi Admission Controller |
% | 3 | Modified Image | Sửa tag sau build (digest mismatch) | Bị chặn |
% | 4 | No Provenance | Image signed nhưng thiếu provenance | Bị chặn |
% | 5 | Untrusted Builder | Provenance từ builder không tin cậy | Bị chặn |
% | 6 | Tag-based Deploy | Deploy dùng tag thay vì digest | Bị chặn (tùy policy) |

\chapter{Triển khai và Đánh giá}\label{ch:implementation}

\section{Môi trường triển khai}

% TODO: Mô tả môi trường
% - Phần cứng: CPU, RAM, OS (WSL2 / Linux)
% - Kubernetes: version, kiểu setup (kubeadm / kind / k3s)
% - Danh sách công cụ và phiên bản:
%   + kubectl, cosign, syft, grype, kyverno
%   + GitHub Actions / Tekton (tùy phương án)
% - Container Runtime: containerd
% - Registry: GHCR / Harbor

\section{Triển khai Pipeline CI/CD}

\subsection{Sample Application}
% TODO: Mô tả app mẫu đơn giản (Go/Python web app)
% - Mục đích: minh họa pipeline, không phải ứng dụng phức tạp
% - Dockerfile, source code

\subsection{Cấu hình CI Pipeline}
% TODO: Trình bày cấu hình pipeline
% Phương án A (GitHub Actions):
% - Workflow YAML: build → sign → generate SBOM → generate provenance → push
% - Sử dụng slsa-github-generator cho SLSA L3 provenance
%
% Phương án B (Tekton - nếu triển khai):
% - Pipeline CRDs
% - Tekton Chains configuration
% - ServiceAccount, RBAC

\subsection{Ký số và Tạo Attestation}
% TODO: 
% - Cosign sign workflow (keyless hoặc key-based)
% - Attach SBOM to registry
% - Verify provenance locally

\section{Triển khai Policy Engine trên Kubernetes}

\subsection{Cài đặt Kyverno}
% TODO: Hướng dẫn cài đặt
% helm install kyverno ...

\subsection{Cấu hình ClusterPolicy}
% TODO: Trình bày YAML configs cho từng policy
% - verify-image-signature
% - verify-slsa-provenance
% - require-image-digest
% - verify-sbom (nếu có)

\section{Kết quả kiểm thử các kịch bản}

\subsection{Kịch bản Happy Path}
% TODO: Screenshots + logs
% Image hợp lệ → deploy thành công

\subsection{Kịch bản tấn công - Attack Scenarios}
% TODO: Cho từng kịch bản:
% - Mô tả thao tác tấn công
% - Output / logs từ Kyverno (denial message)
% - Screenshots minh chứng

% Kịch bản 1: Deploy unsigned image
% Kịch bản 2: Deploy modified image (tag tampering)
% Kịch bản 3: Deploy image without provenance
% Kịch bản 4: Deploy from untrusted builder

\section{Đánh giá hiệu năng}

% TODO: Bảng đo lường
% | Metric | Không có Policy | Có Policy (Kyverno) | Overhead |
% |---|---|---|---|
% | Thời gian deploy Pod | x ms | y ms | z% |
% | CPU usage (Kyverno) | - | x mCPU | - |
% | RAM usage (Kyverno) | - | x Mi | - |
% | Pipeline duration | x min | y min | z% |

\section{Đối chiếu với tiêu chuẩn SLSA Level 3}

% TODO: Bảng đối chiếu chi tiết
% | Yêu cầu SLSA L3 | Mô tả | Đạt/Không đạt | Ghi chú |
% |---|---|---|---|
% | Automated build | Build tự động qua CI | ✓ | GitHub Actions / Tekton |
% | Provenance generation | Tạo provenance in-toto | ✓ | slsa-generator / Chains |
% | Signed provenance | Ký số bởi build platform | ✓ | Cosign keyless |
% | Isolated build | Môi trường cô lập | ✓ | Runner riêng / Pod riêng |
% | Ephemeral environment | Tạm thời, dùng xong hủy | ✓ | ... |
% | Unforgeable provenance | Không thể giả mạo | ✓ | Control plane signing |

\chapter{Kết luận và Hướng phát triển}\label{ch:conclusion}

\section{Kết luận}

% TODO: Tóm tắt những gì đã đạt được
% - Xây dựng thành công pipeline CI/CD đạt chuẩn SLSA Level 3
% - Triển khai cơ chế ký số container image và tạo provenance tự động
% - Quản lý SBOM và quét lỗ hổng bảo mật
% - Triển khai Kubernetes Admission Controller (Kyverno) kiểm soát deploy theo Zero Trust
% - Mô phỏng thành công các kịch bản tấn công, chứng minh hiệu quả hệ thống

\section{Khó khăn gặp phải}

% TODO: Liệt kê các khó khăn
% - Cấu hình keyless signing với OIDC
% - Setup Kubernetes cluster local
% - Tekton Chains configuration (nếu sử dụng)
% - Thiếu tài liệu tiếng Việt
% - Giới hạn tài nguyên phần cứng

\section{Hướng phát triển tương lai}

% TODO: Đề xuất hướng mở rộng
% - Nâng cấp lên SLSA Level 4 (hermetic builds, reproducible builds)
% - Tích hợp DAST (Dynamic Application Security Testing)
% - Runtime security monitoring (Falco)
% - Network Policies cho K8s
% - GitOps deployment với Argo CD / Flux
% - Self-hosted Sigstore infrastructure
% - Multi-cluster policy enforcement
% - Tích hợp với OCI 1.1 referrers API


% Tài liệu tham khảo
\bibliographystyle{plain} 
\bibliography{refs}

\end{document}
